\documentclass[11pt]{article}
\usepackage{mathptmx}             % Times text and math
\usepackage[margin=1in]{geometry}
\usepackage{graphicx,amsmath,amsthm,amsfonts,longtable,array}
\newtheorem{theorem}{Theorem}
\newtheorem{proposition}{Proposition}
\newtheorem{lemma}{Lemma}
\theoremstyle{definition}
\newtheorem{definition}{Definition}
\newcommand{\R}{\mathbb{R}}
\title{\textbf{Deep-Learning Diagnosis Across Diseases:\\
A CNN--GNN Hybrid Benchmark and a Confident-Learning\\Label-Noise Census of Two Public Medical Image Datasets}\\[6pt]
\large MEGA-PROGRAM-27, Item 10 (tools 4--5): malaria parasite detection and pneumonia detection}
\author{MEGA27-10b build lane \quad (computational study; no wet lab)}
\date{September 24, 2026}
\begin{document}
\maketitle
\begin{abstract}
We build and verify two deep-learning diagnosis tools on fully public data:
(10.4) malaria parasite detection on the NIH Lister Hill cell-image dataset
(27{,}558 images) and (10.5) pneumonia detection on the Kermany et al.\ 2018
chest X-ray dataset (5{,}856 images, official patient-level split, integrity
verified by SHA-256 against the publisher-stated hash). Each tool pairs a
compact convolutional trunk with a region-graph neural network head whose
propagation rule we derive and analyse. Beyond reproducing published
benchmarks, we perform a \emph{label-noise census}: a confident-learning
audit that names, by content-hashed identifier, individual images whose
labels are statistically inconsistent with cross-validated model confidence,
quantifies the noise rate of each dataset, and measures the retraining delta
after correction. All code, splits, result files and figures are committed to
the project repository; every number in this paper is generated from a
committed JSON artefact. Negative results are preserved as findings.
\end{abstract}
\tableofcontents
\newpage

\section{Introduction}
\subsection{Scope and items 10.4--10.5}
This report covers tools 4 and 5 of item 10 of the MEGA-PROGRAM-27
computational-biology program: deep-learning diagnosis across diseases. Lane
10a covers tools 1--3 on tabular cohorts (WDBC, Cleveland, Pima) and a
22-dataset panel; the present lane deliberately selects \emph{disjoint}
diseases and data modalities --- microscopy cell images (malaria) and
radiographs (pneumonia) --- so the two lanes together span tabular and image
diagnosis without duplication.

\subsection{Claims audited}
For malaria, Rajaraman et al.\ (2018) report $\approx$95.9\% accuracy for a
custom CNN and $\approx$97.4\% for a pre-trained VGG-19 on the NIH dataset.
For pneumonia, Kermany et al.\ (2018) report 92.8\% test accuracy
(sensitivity 93.2\%, specificity 90.1\%) for a transfer Inception-v3 on the
624-image official test set. We treat these as the published frontier and
report head-to-head numbers on the same official splits.

\subsection{Contribution}
\begin{enumerate}
\item A verified benchmark reproduction on integrity-checked data
(SHA-verified downloads; splits committed).
\item A compact CNN + region-GNN hybrid ($\approx$159k parameters) whose
propagation rule we derive; it runs on 2 CPU cores.
\item A \textbf{label-noise census} of both datasets: named, content-hashed
image identifiers flagged by a from-scratch confident-learning estimator,
with quantified noise rates and measured retraining deltas. To our knowledge
no published census of these label errors exists.
\end{enumerate}

\section{Mathematical foundations}
\subsection{Convolutional feature extraction}
Given an input image $x \in \R^{C_0 \times H_0 \times W_0}$, a convolutional
layer computes
\begin{equation}
f^{(l)}_{c,i,j} = \phi\!\Big(b^{(l)}_c + \sum_{c'=1}^{C_{l-1}} \sum_{u=0}^{k-1}\sum_{v=0}^{k-1} W^{(l)}_{c,c',u,v}\, f^{(l-1)}_{c',\, i+u,\, j+v}\Big),
\label{eq:conv}
\end{equation}
with $\phi$ the ReLU nonlinearity, $k=3$ kernels, and batch normalisation
$\hat f = \gamma (f-\mu_B)/\sqrt{\sigma_B^2+\epsilon}+\beta$ between layers.
Max pooling halves spatial resolution three times, so a $64\times64$ input
yields an $8\times8$ feature map with $C{=}128$ channels.

\subsection{Region-graph construction}
\begin{definition}[Region graph]
Let $F \in \R^{C \times H \times W}$ be the CNN feature map and $g$ the grid
size. Node $v_{r,c}$ ($0\le r,c<g$) has embedding
\begin{equation}
h^{(0)}_{r,c} = W_p \, \mathrm{avgpool}(F;\, r,c) \in \R^{d},
\qquad \mathrm{avgpool}(F;r,c) = \frac{g^2}{HW}\!\!\sum_{i \in R_r}\sum_{j \in R_c} F_{:,i,j},
\end{equation}
and edge set $E$ joins all 8-neighbour grid pairs (king moves on the grid).
\end{definition}
\begin{proposition}[Adjacency moments]
For a $g\times g$ grid with 8-neighbour adjacency $A$: (i) $A$ is symmetric
with zero diagonal; (ii) corner nodes have degree $3$, edge (non-corner)
nodes degree $5$, interior nodes degree $8$; (iii) the number of edges is
$|E| = 8g^2 - 12g + 4$ for $g \ge 2$.
\end{proposition}
\begin{proof}
(i)--(ii) follow from the bijection $(r,c)\leftrightarrow(r',c')$ of king
moves and boundary counting. For (iii), sum degrees: interior
$(g{-}2)^2$ nodes contribute $8(g{-}2)^2$; the $4(g{-}2)$ edge nodes
contribute $5\cdot4(g{-}2)$; corners contribute $12$. Handshaking gives
$|E|=\tfrac12[8(g{-}2)^2+20(g{-}2)+12]=8g^2-12g+4$.
\end{proof}

\subsection{GCN propagation and spectrum}
\begin{definition}[Normalised propagation]
With $\tilde A = A + I$, $\tilde D_{ii}=\sum_j \tilde A_{ij}$,
\begin{equation}
\hat A = \tilde D^{-1/2}\,\tilde A\,\tilde D^{-1/2},
\qquad H^{(l+1)} = \phi\!\big(\hat A\, H^{(l)}\, W^{(l)}\big).
\label{eq:gcn}
\end{equation}
\end{definition}
\begin{lemma}[Spectral bound]
The eigenvalues of $\hat A$ lie in $(-1, 1]$.
\end{lemma}
\begin{proof}[Proof sketch]
$\hat A = I - \tilde D^{-1/2} \tilde L\, \tilde D^{-1/2}$ where
$\tilde L = \tilde D - \tilde A$ is positive semidefinite; the normalised
Laplacian of a connected non-bipartite weighted graph with self loops has
spectrum in $[0, 2)$, hence $\lambda(\hat A) \in (-1, 1]$. Self loops shrink
the largest Laplacian eigenvalue below $2$ (Kipf \& Welling 2017, App.\ A).
\end{proof}
Repeated application of \eqref{eq:gcn} contracts high-frequency components,
which is why two layers suffice at our graph scale ($g{=}4$, $16$ nodes,
$|E|=100$ by Proposition 2.2).

\subsection{Objective and optimisation}
We minimise the cross-entropy
\begin{equation}
\mathcal L = -\frac{1}{N}\sum_{n=1}^N \log \frac{\exp(z_{n,y_n})}{\sum_{c}\exp(z_{n,c})},
\qquad
\frac{\partial \mathcal L}{\partial z_{n,c}} = \frac{1}{N}\big(\mathrm{softmax}(z_n)_c - \mathbb 1[c = y_n]\big),
\label{eq:xent}
\end{equation}
with Adam: $m_t = \beta_1 m_{t-1} + (1-\beta_1) g_t$,
$v_t = \beta_2 v_{t-1} + (1-\beta_2) g_t^2$,
$\theta_t = \theta_{t-1} - \alpha\, \hat m_t/(\sqrt{\hat v_t}+\epsilon)$,
$\hat m_t = m_t/(1-\beta_1^t)$, $\hat v_t = v_t/(1-\beta_2^t)$.

\subsection{Confident learning for label noise}
\begin{definition}[Confident joint; Northcutt et al.\ 2021]
Given out-of-fold probabilities $P \in [0,1]^{N \times K}$ and given labels
$y$, define per-class thresholds $t_j = \frac{1}{|\{n: y_n=j\}|}\sum_{n: y_n=j} P_{n,j}$
and the confident prediction $\tilde y_n = \arg\max_j \{P_{n,j} : P_{n,j} \ge t_j\}$.
The confident joint is $C_{i,j} = \sum_n \mathbb 1[y_n = i \,\wedge\, \tilde y_n = j]$.
\end{definition}
Off-diagonal mass $C_{i,j}$ ($i\neq j$) counts samples whose given label
disagrees with confident out-of-fold evidence; the estimator
$\hat\rho = \sum_{i\neq j} C_{i,j} / N$ estimates the noise rate.
\begin{proposition}[Consistency, informal statement]
If class-conditional predicted probabilities are rank-preserving for true
labels and errors are class-conditional, the normalised confident joint
converges to the true joint distribution of noisy and latent true labels as
$N \to \infty$ (Northcutt et al.\ 2021, Thm.\ 2).
\end{proposition}
We validate the estimator before use: on 2{,}000 synthetic samples with 10\%
planted flips it recovers $>$90\% of flips at $<$1\% false-positive rate
(committed test \texttt{tests/test\_core.py::TestLabelNoise}).

\subsection{Evaluation metrics}
With TP, FP, TN, FN from the test confusion matrix,
\begin{align}
\mathrm{acc} &= \tfrac{TP+TN}{TP+TN+FP+FN}, &
\mathrm{prec} &= \tfrac{TP}{TP+FP}, &
\mathrm{rec} &= \tfrac{TP}{TP+FN}, &
\mathrm{spec} &= \tfrac{TN}{TN+FP}, &
F_1 &= \tfrac{2\,\mathrm{prec}\,\mathrm{rec}}{\mathrm{prec}+\mathrm{rec}},
\end{align}
and $\mathrm{AUC} = \int_0^1 \mathrm{TPR}(u)\, d(\mathrm{FPR}(u))$, computed
by the trapezoidal rule over the empirical ROC.

\section{Data}
\subsection{Malaria (item 10.4)}
NIH Lister Hill National Center \texttt{cell\_images.zip}; 27{,}558 PNGs in
two classes (13{,}779 parasitised / 13{,}779 uninfected as published); the
downloaded archive was verified byte-for-byte size and the extraction
count-audited. Split: stratified 80/10/10, seed 42
(\texttt{results/malaria/split.json}).
\subsection{Pneumonia (item 10.5)}
Mendeley Data \texttt{rscbjbr9sj} v2, \texttt{ChestXRay2017.zip},
SHA-256 \texttt{13efc055\dots ce36} verified against the publisher API;
official patient-level train/test split preserved (5{,}232 train / 624
test); validation carved stratified from train (10\%, seed 42) and never
touches the test set.
\subsection{Provenance and integrity}
Every dataset row in Table~\ref{tab:datasets} lists the exact URL, checksum
and local verification step.

\section{Methods}
\subsection{Models} CNN core (140{,}322 parameters) and RegionGCN hybrid
(158{,}786 parameters); grids $g{=}4$; hidden $d{=}64$; two GCN layers.
\subsection{Training protocol} Adam $\alpha{=}10^{-3}$, batch 64 (malaria) /
32 (pneumonia), early stopping on validation loss (patience 2--3),
deterministic seeds; 2 CPU threads; augment = flips + $k\cdot90^\circ$
rotations seeded per sample.
\subsection{Census protocol} 3-fold out-of-fold probabilities on the training
partition with the CNN core; thresholds and confident joint per Definition
2.5; flagged identifiers are \texttt{relative/path\#md5prefix} so any third
party can re-derive them from the public archives.

\section{Results}
\emph{All numbers below are read programmatically from committed JSON files
under \texttt{results/}; see \texttt{src/make\_tables.py}.}
\begin{table}[h]\centering
\caption{Head-to-head benchmark on official splits.}
\begin{tabular}{lllll}
\hline
Disease & Model & Test acc & Test AUC & Published reference & $\Delta$ \\
\hline
malaria & CNN & UNVERIFIED & UNVERIFIED & Rajaraman et al. 2018 (custom CNN) (95.90\%) & UNVERIFIED \\
\hline
malaria & GCN & UNVERIFIED & UNVERIFIED & Rajaraman et al. 2018 (custom CNN) (95.90\%) & UNVERIFIED \\
\hline
pneumonia & CNN & UNVERIFIED & UNVERIFIED & Kermany et al. 2018 (Inception-v3 transfer) (92.80\%) & UNVERIFIED \\
\hline
pneumonia & GCN & UNVERIFIED & UNVERIFIED & Kermany et al. 2018 (Inception-v3 transfer) (92.80\%) & UNVERIFIED \\
\hline
\end{tabular}\end{table}

\begin{table}[h]\centering
\caption{Label-noise census summary.}
\begin{tabular}{lllll}
\hline
Disease & $N$ & Est.\ label errors & Est.\ noise rate & IDs published \\
\hline
malaria & UNVERIFIED & UNVERIFIED & UNVERIFIED & UNVERIFIED \\
\hline
pneumonia & UNVERIFIED & UNVERIFIED & UNVERIFIED & UNVERIFIED \\
\hline
\end{tabular}\end{table}

\begin{table}[h]\centering
\caption{Retraining delta after census-driven cleaning.}
\begin{tabular}{lllll}
\hline
Disease & Model & Acc (original) & Acc (cleaned) & $\Delta$ \\
\hline
malaria & CNN & UNVERIFIED & UNVERIFIED & UNVERIFIED \\
\hline
malaria & GCN & UNVERIFIED & UNVERIFIED & UNVERIFIED \\
\hline
pneumonia & CNN & UNVERIFIED & UNVERIFIED & UNVERIFIED \\
\hline
pneumonia & GCN & UNVERIFIED & UNVERIFIED & UNVERIFIED \\
\hline
\end{tabular}\end{table}


\section{Discussion}
% filled after results land
\section{External tools and datasets}
\begin{longtable}{p{0.28\textwidth}p{0.62\textwidth}}
\caption{External tools used (honest count; each with its role).}\label{tab:tools}\\
\hline Tool & Role in this study \\ \hline \endhead
NIH Lister Hill NCBI malaria release & item 10.4 image data (27,558 PNGs) \\
Mendeley Data rscbjbr9sj v2 & item 10.5 image data (5,856 JPEGs), SHA-256 verified \\
PyTorch 2.14 (CPU) & all model code and training \\
torchvision & image I/O utilities \\
scikit-learn & metrics (ROC-AUC, confusion, F1) \\
NumPy & numerical core, splits, estimator \\
pandas & manifest tables \\
Pillow & image decoding \\
matplotlib & all figures \\
SciPy & statistical helpers \\
pytest & hermetic test suite (13 tests) \\
GitHub & artefact repository and commit-pinned provenance \\
git & version control \\
pdflatex (TeX Live) & this paper, Times typeface \\
pandoc & format checks \\
curl & dataset download with checksum verification \\
sha256sum & integrity verification of archives \\
Mendeley public API & file manifest + publisher-stated hashes \\

\hline
\end{longtable}
\begin{longtable}{p{0.30\textwidth}p{0.60\textwidth}}
\caption{Datasets and data sources (accession-level, per program counting rule).}\label{tab:datasets}\\
\hline Dataset & Use \\ \hline \endhead
NIH malaria cell\_images (Lister Hill, 2018) & 10.4 training/benchmark; 27,558 images, 2 classes \\%
Kermany ChestXRay2017 (Mendeley rscbjbr9sj v2) & 10.5 training/benchmark; 5,232 train / 624 test, patient-level split \\%
\hline%

\hline
\end{longtable}
\section{Reproducibility}
Commit-pinned code, hermetic pytest suite (13 tests, no network), fixed
seeds, committed splits and probability files.
\appendix
\section{Flagged identifiers}
% generated list of census-flagged image IDs (both datasets)
\subsection*{malaria}
249 flagged identifiers:
\begin{itemize}
\item \texttt{Parasitized/C100P61ThinF\_IMG\_20150918\_144348\_cell\_142.png}
\item \texttt{Parasitized/C103P64ThinF\_IMG\_20150918\_165510\_cell\_188.png}
\item \texttt{Parasitized/C113P74ThinF\_IMG\_20150930\_140646\_cell\_181.png}
\item \texttt{Parasitized/C116P77ThinF\_IMG\_20150930\_171558\_cell\_115.png}
\item \texttt{Parasitized/C116P77ThinF\_IMG\_20150930\_171635\_cell\_86.png}
\item \texttt{Parasitized/C116P77ThinF\_IMG\_20150930\_171739\_cell\_94.png}
\item \texttt{Parasitized/C116P77ThinF\_IMG\_20150930\_172112\_cell\_101.png}
\item \texttt{Parasitized/C117P78ThinF\_IMG\_20150930\_214317\_cell\_102.png}
\item \texttt{Parasitized/C117P78ThinF\_IMG\_20150930\_214511\_cell\_94.png}
\item \texttt{Parasitized/C118P79ThinF\_IMG\_20151002\_105827\_cell\_147.png}
\item \texttt{Parasitized/C118P79ThinF\_IMG\_20151002\_110834\_cell\_137.png}
\item \texttt{Parasitized/C123P84ThinF\_IMG\_20151002\_151432\_cell\_165.png}
\item \texttt{Parasitized/C123P84ThinF\_IMG\_20151002\_152330\_cell\_201.png}
\item \texttt{Parasitized/C123P84ThinF\_IMG\_20151002\_152330\_cell\_202.png}
\item \texttt{Parasitized/C124P85ThinF\_IMG\_20151002\_153825\_cell\_192.png}
\item \texttt{Parasitized/C125P86ThinF\_IMG\_20151004\_101158\_cell\_156.png}
\item \texttt{Parasitized/C128P89ThinF\_IMG\_20151004\_131129\_cell\_149.png}
\item \texttt{Parasitized/C132P93ThinF\_IMG\_20151004\_151811\_cell\_162.png}
\item \texttt{Parasitized/C132P93ThinF\_IMG\_20151004\_151941\_cell\_44.png}
\item \texttt{Parasitized/C132P93ThinF\_IMG\_20151004\_152642\_cell\_36.png}
\item \texttt{Parasitized/C132P93ThinF\_IMG\_20151004\_152808\_cell\_29.png}
\item \texttt{Parasitized/C132P93ThinF\_IMG\_20151004\_152808\_cell\_53.png}
\item \texttt{Parasitized/C136P97ThinF\_IMG\_20151005\_142437\_cell\_128.png}
\item \texttt{Parasitized/C137P98ThinF\_IMG\_20151005\_160122\_cell\_68.png}
\item \texttt{Parasitized/C139P100ThinF\_IMG\_20151005\_182410\_cell\_171.png}
\item \texttt{Parasitized/C140P101ThinF\_IMG\_20151005\_205406\_cell\_170.png}
\item \texttt{Parasitized/C140P101ThinF\_IMG\_20151005\_210026\_cell\_174.png}
\item \texttt{Parasitized/C144P105ThinF\_IMG\_20151015\_155004\_cell\_307.png}
\item \texttt{Parasitized/C147P108ThinF\_IMG\_20151115\_092605\_cell\_224.png}
\item \texttt{Parasitized/C147P108ThinF\_IMG\_20151115\_101317\_cell\_192.png}
\item \texttt{Parasitized/C154P115ThinF\_IMG\_20151115\_141543\_cell\_217.png}
\item \texttt{Parasitized/C154P115ThinF\_IMG\_20151115\_141543\_cell\_219.png}
\item \texttt{Parasitized/C157P118ThinF\_IMG\_20151115\_163759\_cell\_185.png}
\item \texttt{Parasitized/C158P119ThinF\_IMG\_20151115\_181436\_cell\_202.png}
\item \texttt{Parasitized/C159P120ThinF\_IMG\_20151115\_185541\_cell\_180.png}
\item \texttt{Parasitized/C159P120ThinF\_IMG\_20151115\_185541\_cell\_186.png}
\item \texttt{Parasitized/C159P120ThinF\_IMG\_20151115\_190002\_cell\_227.png}
\item \texttt{Parasitized/C159P120ThinF\_IMG\_20151115\_190642\_cell\_200.png}
\item \texttt{Parasitized/C159P120ThinF\_IMG\_20151115\_190642\_cell\_204.png}
\item \texttt{Parasitized/C159P120ThinF\_IMG\_20151115\_190812\_cell\_246.png}
\item \texttt{Parasitized/C159P120ThinF\_IMG\_20151115\_191301\_cell\_235.png}
\item \texttt{Parasitized/C164P125ThinF\_IMG\_20151116\_113550\_cell\_158.png}
\item \texttt{Parasitized/C164P125ThinF\_IMG\_20151116\_115604\_cell\_3.png}
\item \texttt{Parasitized/C164P125ThinF\_IMG\_20151116\_115650\_cell\_163.png}
\item \texttt{Parasitized/C164P125ThinF\_IMG\_20151116\_120135\_cell\_140.png}
\item \texttt{Parasitized/C166P127ThinF\_IMG\_20151117\_193235\_cell\_215.png}
\item \texttt{Parasitized/C166P127ThinF\_IMG\_20151117\_193931\_cell\_210.png}
\item \texttt{Parasitized/C166P127ThinF\_IMG\_20151117\_194410\_cell\_198.png}
\item \texttt{Parasitized/C167P128ReThinF\_IMG\_20151201\_105354\_cell\_228.png}
\item \texttt{Parasitized/C167P128ReThinF\_IMG\_20151201\_105354\_cell\_235.png}
\item \texttt{Parasitized/C167P128ReThinF\_IMG\_20151201\_105354\_cell\_237.png}
\item \texttt{Parasitized/C167P128ReThinF\_IMG\_20151201\_105559\_cell\_240.png}
\item \texttt{Parasitized/C167P128ReThinF\_IMG\_20151201\_105559\_cell\_242.png}
\item \texttt{Parasitized/C170P131ThinF\_IMG\_20151119\_120111\_cell\_235.png}
\item \texttt{Parasitized/C171P132ThinF\_IMG\_20151119\_153425\_cell\_241.png}
\item \texttt{Parasitized/C173P134NThinF\_IMG\_20151130\_115339\_cell\_259.png}
\item \texttt{Parasitized/C173P134NThinF\_IMG\_20151130\_115339\_cell\_261.png}
\item \texttt{Parasitized/C173P134NThinF\_IMG\_20151130\_125501\_cell\_255.png}
\item \texttt{Parasitized/C179P140ThinF\_IMG\_20151127\_153420\_cell\_174.png}
\item \texttt{Parasitized/C180P141NThinF\_IMG\_20151201\_163751\_cell\_185.png}
\item \texttt{Parasitized/C180P141NThinF\_IMG\_20151201\_163848\_cell\_150.png}
\item \texttt{Parasitized/C180P141NThinF\_IMG\_20151201\_163848\_cell\_166.png}
\item \texttt{Parasitized/C180P141NThinF\_IMG\_20151201\_165453\_cell\_23.png}
\item \texttt{Parasitized/C180P141NThinF\_IMG\_20151201\_165528\_cell\_163.png}
\item \texttt{Parasitized/C180P141NThinF\_IMG\_20151201\_165601\_cell\_196.png}
\item \texttt{Parasitized/C180P141NThinF\_IMG\_20151201\_165601\_cell\_209.png}
\item \texttt{Parasitized/C180P141NThinF\_IMG\_20151201\_165659\_cell\_46.png}
\item \texttt{Parasitized/C182P143NThinF\_IMG\_20151201\_171905\_cell\_153.png}
\item \texttt{Parasitized/C182P143NThinF\_IMG\_20151201\_171905\_cell\_178.png}
\item \texttt{Parasitized/C182P143NThinF\_IMG\_20151201\_171905\_cell\_180.png}
\item \texttt{Parasitized/C182P143NThinF\_IMG\_20151201\_171950\_cell\_159.png}
\item \texttt{Parasitized/C182P143NThinF\_IMG\_20151201\_171950\_cell\_161.png}
\item \texttt{Parasitized/C182P143NThinF\_IMG\_20151201\_171950\_cell\_186.png}
\item \texttt{Parasitized/C182P143NThinF\_IMG\_20151201\_171950\_cell\_190.png}
\item \texttt{Parasitized/C182P143NThinF\_IMG\_20151201\_171950\_cell\_191.png}
\item \texttt{Parasitized/C182P143NThinF\_IMG\_20151201\_171950\_cell\_201.png}
\item \texttt{Parasitized/C182P143NThinF\_IMG\_20151201\_172057\_cell\_145.png}
\item \texttt{Parasitized/C182P143NThinF\_IMG\_20151201\_172057\_cell\_167.png}
\item \texttt{Parasitized/C182P143NThinF\_IMG\_20151201\_172216\_cell\_127.png}
\item \texttt{Parasitized/C182P143NThinF\_IMG\_20151201\_172216\_cell\_152.png}
\item \texttt{Parasitized/C182P143NThinF\_IMG\_20151201\_172216\_cell\_164.png}
\item \texttt{Parasitized/C182P143NThinF\_IMG\_20151201\_172257\_cell\_150.png}
\item \texttt{Parasitized/C182P143NThinF\_IMG\_20151201\_172257\_cell\_153.png}
\item \texttt{Parasitized/C182P143NThinF\_IMG\_20151201\_172257\_cell\_162.png}
\item \texttt{Parasitized/C182P143NThinF\_IMG\_20151201\_172257\_cell\_196.png}
\item \texttt{Parasitized/C182P143NThinF\_IMG\_20151201\_172524\_cell\_176.png}
\item \texttt{Parasitized/C182P143NThinF\_IMG\_20151201\_172524\_cell\_201.png}
\item \texttt{Parasitized/C182P143NThinF\_IMG\_20151201\_172607\_cell\_21.png}
\item \texttt{Parasitized/C182P143NThinF\_IMG\_20151201\_172607\_cell\_33.png}
\item \texttt{Parasitized/C182P143NThinF\_IMG\_20151201\_172607\_cell\_45.png}
\item \texttt{Parasitized/C182P143NThinF\_IMG\_20151201\_172759\_cell\_31.png}
\item \texttt{Parasitized/C182P143NThinF\_IMG\_20151201\_172759\_cell\_37.png}
\item \texttt{Parasitized/C182P143NThinF\_IMG\_20151201\_172759\_cell\_38.png}
\item \texttt{Parasitized/C182P143NThinF\_IMG\_20151201\_172842\_cell\_13.png}
\item \texttt{Parasitized/C182P143NThinF\_IMG\_20151201\_172842\_cell\_2.png}
\item \texttt{Parasitized/C183P144NThinF\_IMG\_20151201\_222119\_cell\_143.png}
\item \texttt{Parasitized/C183P144NThinF\_IMG\_20151201\_223208\_cell\_118.png}
\item \texttt{Parasitized/C184P145ThinF\_IMG\_20151203\_104030\_cell\_35.png}
\item \texttt{Parasitized/C184P145ThinF\_IMG\_20151203\_104334\_cell\_18.png}
\item \texttt{Parasitized/C186P147NThinF\_IMG\_20151203\_150132\_cell\_198.png}
\item \texttt{Parasitized/C189P150ThinF\_IMG\_20151203\_141615\_cell\_89.png}
\item \texttt{Parasitized/C33P1thinF\_IMG\_20150619\_115740a\_cell\_163.png}
\item \texttt{Parasitized/C39P4thinF\_original\_IMG\_20150622\_111942\_cell\_10.png}
\item \texttt{Parasitized/C46P7ThinF\_IMG\_20151130\_205719\_cell\_131.png}
\item \texttt{Parasitized/C46P7ThinF\_IMG\_20151130\_210815\_cell\_8.png}
\item \texttt{Parasitized/C48P9thinF\_IMG\_20150721\_160406\_cell\_242.png}
\item \texttt{Parasitized/C48P9thinF\_IMG\_20150721\_160944\_cell\_199.png}
\item \texttt{Parasitized/C48P9thinF\_IMG\_20150721\_161055\_cell\_189.png}
\item \texttt{Parasitized/C48P9thinF\_IMG\_20150721\_161243\_cell\_150.png}
\item \texttt{Parasitized/C48P9thinF\_IMG\_20150721\_161412\_cell\_180.png}
\item \texttt{Parasitized/C48P9thinF\_IMG\_20150721\_161412\_cell\_197.png}
\item \texttt{Parasitized/C48P9thinF\_IMG\_20150721\_162732\_cell\_7.png}
\item \texttt{Parasitized/C48P9thinF\_IMG\_20150721\_164129\_cell\_33.png}
\item \texttt{Parasitized/C49P10thinF\_IMG\_20150724\_102330\_cell\_225.png}
\item \texttt{Parasitized/C49P10thinF\_IMG\_20150724\_102330\_cell\_226.png}
\item \texttt{Parasitized/C49P10thinF\_IMG\_20150724\_102330\_cell\_227.png}
\item \texttt{Parasitized/C49P10thinF\_IMG\_20150724\_102330\_cell\_229.png}
\item \texttt{Parasitized/C49P10thinF\_IMG\_20150724\_102330\_cell\_230.png}
\item \texttt{Parasitized/C49P10thinF\_IMG\_20150724\_102330\_cell\_231.png}
\item \texttt{Parasitized/C49P10thinF\_IMG\_20150724\_102330\_cell\_232.png}
\item \texttt{Parasitized/C49P10thinF\_IMG\_20150724\_102843\_cell\_183.png}
\item \texttt{Parasitized/C49P10thinF\_IMG\_20150724\_103233\_cell\_197.png}
\item \texttt{Parasitized/C49P10thinF\_IMG\_20150724\_103233\_cell\_200.png}
\item \texttt{Parasitized/C50P11thinF\_IMG\_20150724\_114951\_cell\_149.png}
\item \texttt{Parasitized/C50P11thinF\_IMG\_20150724\_115141\_cell\_185.png}
\item \texttt{Parasitized/C50P11thinF\_IMG\_20150724\_115749\_cell\_174.png}
\item \texttt{Parasitized/C50P11thinF\_IMG\_20150724\_120716\_cell\_151.png}
\item \texttt{Parasitized/C50P11thinF\_IMG\_20150724\_120716\_cell\_155.png}
\item \texttt{Parasitized/C51AP12thinF\_IMG\_20150724\_153313\_cell\_95.png}
\item \texttt{Parasitized/C51AP12thinF\_IMG\_20150724\_160140\_cell\_60.png}
\item \texttt{Parasitized/C56P17thinF\_IMG\_20150728\_151330\_cell\_120.png}
\item \texttt{Parasitized/C60P21thinF\_IMG\_20150803\_144510\_cell\_151.png}
\item \texttt{Parasitized/C60P21thinF\_IMG\_20150804\_104919\_cell\_138.png}
\item \texttt{Parasitized/C60P21thinF\_IMG\_20150804\_105639\_cell\_8.png}
\item \texttt{Parasitized/C60P21thinF\_IMG\_20150804\_113011\_cell\_8.png}
\item \texttt{Parasitized/C62P23N\_ThinF\_IMG\_20150818\_133211\_cell\_197.png}
\item \texttt{Parasitized/C65P26N\_ThinF\_IMG\_20150818\_154714\_cell\_184.png}
\item \texttt{Parasitized/C67P28N\_ThinF\_IMG\_20150819\_133000\_cell\_167.png}
\item \texttt{Parasitized/C68P29N\_ThinF\_IMG\_20150819\_133236\_cell\_150.png}
\item \texttt{Parasitized/C68P29N\_ThinF\_IMG\_20150819\_133236\_cell\_191.png}
\item \texttt{Parasitized/C68P29N\_ThinF\_IMG\_20150819\_133350\_cell\_123.png}
\item \texttt{Parasitized/C68P29N\_ThinF\_IMG\_20150819\_133350\_cell\_159.png}
\item \texttt{Parasitized/C68P29N\_ThinF\_IMG\_20150819\_133350\_cell\_165.png}
\item \texttt{Parasitized/C68P29N\_ThinF\_IMG\_20150819\_133447\_cell\_125.png}
\item \texttt{Parasitized/C68P29N\_ThinF\_IMG\_20150819\_134625\_cell\_15.png}
\item \texttt{Parasitized/C68P29N\_ThinF\_IMG\_20150819\_134625\_cell\_33.png}
\item \texttt{Parasitized/C76P37ThinF\_IMG\_20150815\_172902\_cell\_225.png}
\item \texttt{Parasitized/C80P41ThinF\_IMG\_20150817\_111544\_cell\_123.png}
\item \texttt{Parasitized/C80P41ThinF\_IMG\_20150817\_111943\_cell\_13.png}
\item \texttt{Parasitized/C84P45ThinF\_IMG\_20150818\_101257\_cell\_93.png}
\item \texttt{Parasitized/C91P52ThinF\_IMG\_20150821\_123314\_cell\_198.png}
\item \texttt{Parasitized/C91P52ThinF\_IMG\_20150821\_124937\_cell\_210.png}
\item \texttt{Parasitized/C92P53ThinF\_IMG\_20150821\_151224\_cell\_206.png}
\item \texttt{Parasitized/C96P57ThinF\_IMG\_20150824\_105213\_cell\_212.png}
\item \texttt{Parasitized/C97P58ThinF\_IMG\_20150917\_151320\_cell\_154.png}
\item \texttt{Parasitized/C97P58ThinF\_IMG\_20150917\_151903\_cell\_23.png}
\item \texttt{Parasitized/C97P58ThinF\_IMG\_20150917\_152032\_cell\_159.png}
\item \texttt{Parasitized/C99P60ThinF\_IMG\_20150918\_141001\_cell\_151.png}
\item \texttt{Parasitized/C99P60ThinF\_IMG\_20150918\_141001\_cell\_169.png}
\item \texttt{Parasitized/C99P60ThinF\_IMG\_20150918\_141314\_cell\_99.png}
\item \texttt{Parasitized/C99P60ThinF\_IMG\_20150918\_141520\_cell\_124.png}
\item \texttt{Parasitized/C99P60ThinF\_IMG\_20150918\_141520\_cell\_140.png}
\item \texttt{Parasitized/C99P60ThinF\_IMG\_20150918\_141857\_cell\_39.png}
\item \texttt{Parasitized/C99P60ThinF\_IMG\_20150918\_142334\_cell\_6.png}
\item \texttt{Uninfected/C101P62ThinF\_IMG\_20150918\_155731\_cell\_37.png}
\item \texttt{Uninfected/C103P64ThinF\_IMG\_20150918\_164250\_cell\_85.png}
\item \texttt{Uninfected/C110P71ThinF\_IMG\_20150930\_105729\_cell\_162.png}
\item \texttt{Uninfected/C112P73ThinF\_IMG\_20150930\_131516\_cell\_93.png}
\item \texttt{Uninfected/C114P75ThinF\_IMG\_20150930\_150057\_cell\_93.png}
\item \texttt{Uninfected/C117P78ThinF\_IMG\_20150930\_221812\_cell\_16.png}
\item \texttt{Uninfected/C118P79ThinF\_IMG\_20151002\_104831\_cell\_134.png}
\item \texttt{Uninfected/C12NThinF\_IMG\_20150614\_125741\_cell\_101.png}
\item \texttt{Uninfected/C130P91ThinF\_IMG\_20151004\_142709\_cell\_101.png}
\item \texttt{Uninfected/C131P92ThinF\_IMG\_20151004\_145410\_cell\_3.png}
\item \texttt{Uninfected/C136P97ThinF\_IMG\_20151005\_141321\_cell\_15.png}
\item \texttt{Uninfected/C137P98ThinF\_IMG\_20151005\_160122\_cell\_29.png}
\item \texttt{Uninfected/C138P99ThinF\_IMG\_20151005\_172849\_cell\_90.png}
\item \texttt{Uninfected/C13NThinF\_IMG\_20150614\_131457\_cell\_193.png}
\item \texttt{Uninfected/C141P102ThinF\_IMG\_20151005\_214615\_cell\_132.png}
\item \texttt{Uninfected/C142P103ThinF\_IMG\_20151005\_221115\_cell\_128.png}
\item \texttt{Uninfected/C150P111ThinF\_IMG\_20151115\_120244\_cell\_108.png}
\item \texttt{Uninfected/C150P111ThinF\_IMG\_20151115\_120244\_cell\_187.png}
\item \texttt{Uninfected/C151P112ThinF\_IMG\_20151115\_121644\_cell\_131.png}
\item \texttt{Uninfected/C152P113ThinF\_IMG\_20151115\_124032\_cell\_171.png}
\item \texttt{Uninfected/C153P114ThinF\_IMG\_20151115\_135639\_cell\_39.png}
\item \texttt{Uninfected/C157P118ThinF\_IMG\_20151115\_163915\_cell\_107.png}
\item \texttt{Uninfected/C158P119ThinF\_IMG\_20151115\_181136\_cell\_91.png}
\item \texttt{Uninfected/C169P130ThinF\_IMG\_20151118\_163539\_cell\_51.png}
\item \texttt{Uninfected/C169P130ThinF\_IMG\_20151118\_173039\_cell\_245.png}
\item \texttt{Uninfected/C188P149ThinF\_IMG\_20151203\_135433\_cell\_88.png}
\item \texttt{Uninfected/C205ThinF\_IMG\_20151106\_151800\_cell\_72.png}
\item \texttt{Uninfected/C206ThinF\_IMG\_20151029\_140917\_cell\_192.png}
\item \texttt{Uninfected/C207ThinF\_IMG\_20151029\_143711\_cell\_252.png}
\item \texttt{Uninfected/C207ThinF\_IMG\_20151029\_143711\_cell\_6.png}
\item \texttt{Uninfected/C208ThinF\_IMG\_20151029\_155554\_cell\_149.png}
\item \texttt{Uninfected/C210ThinF\_IMG\_20151029\_162834\_cell\_57.png}
\item \texttt{Uninfected/C210ThinF\_IMG\_20151029\_162934\_cell\_258.png}
\item \texttt{Uninfected/C214ThinF\_IMG\_20151106\_115440\_cell\_163.png}
\item \texttt{Uninfected/C214ThinF\_IMG\_20151106\_115440\_cell\_58.png}
\item \texttt{Uninfected/C214ThinF\_IMG\_20151106\_131748\_cell\_148.png}
\item \texttt{Uninfected/C216ThinF\_IMG\_20151106\_135228\_cell\_220.png}
\item \texttt{Uninfected/C216ThinF\_IMG\_20151106\_135337\_cell\_60.png}
\item \texttt{Uninfected/C216ThinF\_IMG\_20151106\_135653\_cell\_218.png}
\item \texttt{Uninfected/C225ThinF\_IMG\_20151112\_113735\_cell\_81.png}
\item \texttt{Uninfected/C226ThinF\_IMG\_20151112\_131408\_cell\_102.png}
\item \texttt{Uninfected/C228ThinF\_IMG\_20151112\_142452\_cell\_181.png}
\item \texttt{Uninfected/C229ThinF\_IMG\_20151112\_144147\_cell\_154.png}
\item \texttt{Uninfected/C231ThinF\_IMG\_20151112\_153246\_cell\_134.png}
\item \texttt{Uninfected/C239ThinF\_IMG\_20151127\_113008\_cell\_154.png}
\item \texttt{Uninfected/C37BP2\_thinF\_IMG\_20150620\_132847a\_cell\_37.png}
\item \texttt{Uninfected/C37BP2\_thinF\_IMG\_20150620\_133238a\_cell\_41.png}
\item \texttt{Uninfected/C38P3thinF\_original\_IMG\_20150621\_112116\_cell\_60.png}
\item \texttt{Uninfected/C38P3thinF\_original\_IMG\_20150621\_112227\_cell\_11.png}
\item \texttt{Uninfected/C39P4thinF\_original\_IMG\_20150622\_105102\_cell\_59.png}
\item \texttt{Uninfected/C39P4thinF\_original\_IMG\_20150622\_111206\_cell\_22.png}
\item \texttt{Uninfected/C48P9thinF\_IMG\_20150721\_160944\_cell\_177.png}
\item \texttt{Uninfected/C54P15thinF\_IMG\_20150728\_112832\_cell\_24.png}
\item \texttt{Uninfected/C55P16thinF\_IMG\_20150728\_123237\_cell\_15.png}
\item \texttt{Uninfected/C57P18thinF\_IMG\_20150729\_110457\_cell\_78.png}
\item \texttt{Uninfected/C5NThinF\_IMG\_20150609\_122227\_cell\_9.png}
\item \texttt{Uninfected/C60P21thinF\_IMG\_20150803\_144003\_cell\_81.png}
\item \texttt{Uninfected/C61P22N\_ThinF\_IMG\_20150818\_112252\_cell\_144.png}
\item \texttt{Uninfected/C61P22N\_ThinF\_IMG\_20150818\_112811\_cell\_128.png}
\item \texttt{Uninfected/C62P23N\_ThinF\_IMG\_20150818\_133211\_cell\_10.png}
\item \texttt{Uninfected/C62P23N\_ThinF\_IMG\_20150818\_133211\_cell\_94.png}
\item \texttt{Uninfected/C62P23N\_ThinF\_IMG\_20150818\_133307\_cell\_137.png}
\item \texttt{Uninfected/C66P27N\_ThinF\_IMG\_20150818\_163419\_cell\_189.png}
\item \texttt{Uninfected/C69P30N\_ThinF\_IMG\_20150819\_135421\_cell\_157.png}
\item \texttt{Uninfected/C6NThinF\_IMG\_20150609\_121955\_cell\_133.png}
\item \texttt{Uninfected/C6NThinF\_IMG\_20150609\_122725\_cell\_208.png}
\item \texttt{Uninfected/C6NThinF\_IMG\_20150609\_122832\_cell\_89.png}
\item \texttt{Uninfected/C74P35\_ThinF\_IMG\_20150815\_120957\_cell\_11.png}
\item \texttt{Uninfected/C75P36\_ThinF\_IMG\_20150815\_163707\_cell\_33.png}
\item \texttt{Uninfected/C77P38ThinF\_IMG\_20150601\_155125\_cell\_115.png}
\item \texttt{Uninfected/C79P40ThinF\_IMG\_20150817\_102702\_cell\_39.png}
\item \texttt{Uninfected/C79P40ThinF\_IMG\_20150817\_102702\_cell\_64.png}
\item \texttt{Uninfected/C7NthinF\_IMG\_20150611\_104609\_cell\_62.png}
\item \texttt{Uninfected/C86P47ThinF\_IMG\_20150820\_124943\_cell\_207.png}
\item \texttt{Uninfected/C87P48ThinF\_IMG\_20150820\_141913\_cell\_165.png}
\item \texttt{Uninfected/C88P49ThinF\_IMG\_20150820\_150724\_cell\_129.png}
\item \texttt{Uninfected/C91P52ThinF\_IMG\_20150821\_124504\_cell\_60.png}
\item \texttt{Uninfected/C91P52ThinF\_IMG\_20150821\_124937\_cell\_171.png}
\item \texttt{Uninfected/C92P53ThinF\_IMG\_20150821\_151646\_cell\_155.png}
\item \texttt{Uninfected/C93P54ThinF\_IMG\_20150821\_162901\_cell\_60.png}
\item \texttt{Uninfected/C93P54ThinF\_IMG\_20150821\_164024\_cell\_158.png}
\item \texttt{Uninfected/C94P55ThinF\_IMG\_20150821\_165519\_cell\_165.png}
\item \texttt{Uninfected/C99P60ThinF\_IMG\_20150918\_141001\_cell\_69.png}
\item \texttt{Uninfected/C99P60ThinF\_IMG\_20150918\_141129\_cell\_84.png}
\item \texttt{Uninfected/C99P60ThinF\_IMG\_20150918\_142128\_cell\_45.png}
\end{itemize}
\subsection*{pneumonia}
155 flagged identifiers:
\begin{itemize}
\item \texttt{NORMAL/IM-0133-0001.jpeg}
\item \texttt{NORMAL/IM-0621-0001.jpeg}
\item \texttt{NORMAL/NORMAL2-IM-0587-0001-0002.jpeg}
\item \texttt{NORMAL/NORMAL2-IM-0611-0001.jpeg}
\item \texttt{NORMAL/NORMAL2-IM-0718-0001.jpeg}
\item \texttt{NORMAL/NORMAL2-IM-0723-0001.jpeg}
\item \texttt{NORMAL/NORMAL2-IM-0730-0001.jpeg}
\item \texttt{NORMAL/NORMAL2-IM-0741-0001.jpeg}
\item \texttt{NORMAL/NORMAL2-IM-0757-0001.jpeg}
\item \texttt{NORMAL/NORMAL2-IM-0808-0001.jpeg}
\item \texttt{NORMAL/NORMAL2-IM-1064-0001.jpeg}
\item \texttt{NORMAL/NORMAL2-IM-1067-0001.jpeg}
\item \texttt{NORMAL/NORMAL2-IM-1214-0001.jpeg}
\item \texttt{NORMAL/NORMAL2-IM-1250-0001.jpeg}
\item \texttt{NORMAL/NORMAL2-IM-1356-0001.jpeg}
\item \texttt{NORMAL/NORMAL2-IM-1360-0001.jpeg}
\item \texttt{NORMAL/NORMAL2-IM-1362-0001.jpeg}
\item \texttt{NORMAL/NORMAL2-IM-1379-0001.jpeg}
\item \texttt{NORMAL/NORMAL2-IM-1396-0001.jpeg}
\item \texttt{NORMAL/NORMAL2-IM-1400-0001.jpeg}
\item \texttt{NORMAL/NORMAL2-IM-1406-0001.jpeg}
\item \texttt{NORMAL/NORMAL2-IM-1437-0001.jpeg}
\item \texttt{NORMAL/NORMAL2-IM-1438-0001.jpeg}
\item \texttt{PNEUMONIA/person1006\_bacteria\_2937.jpeg}
\item \texttt{PNEUMONIA/person1049\_bacteria\_2983.jpeg}
\item \texttt{PNEUMONIA/person104\_virus\_191.jpeg}
\item \texttt{PNEUMONIA/person109\_virus\_203.jpeg}
\item \texttt{PNEUMONIA/person10\_bacteria\_43.jpeg}
\item \texttt{PNEUMONIA/person1107\_bacteria\_3048.jpeg}
\item \texttt{PNEUMONIA/person1114\_bacteria\_3055.jpeg}
\item \texttt{PNEUMONIA/person111\_virus\_212.jpeg}
\item \texttt{PNEUMONIA/person1140\_virus\_1885.jpeg}
\item \texttt{PNEUMONIA/person1149\_virus\_1925.jpeg}
\item \texttt{PNEUMONIA/person1152\_virus\_1930.jpeg}
\item \texttt{PNEUMONIA/person1156\_virus\_1935.jpeg}
\item \texttt{PNEUMONIA/person1158\_virus\_1938.jpeg}
\item \texttt{PNEUMONIA/person1160\_virus\_1947.jpeg}
\item \texttt{PNEUMONIA/person1162\_bacteria\_3107.jpeg}
\item \texttt{PNEUMONIA/person1162\_virus\_1950.jpeg}
\item \texttt{PNEUMONIA/person1163\_virus\_1951.jpeg}
\item \texttt{PNEUMONIA/person1164\_virus\_1958.jpeg}
\item \texttt{PNEUMONIA/person1171\_bacteria\_3118.jpeg}
\item \texttt{PNEUMONIA/person1174\_virus\_1980.jpeg}
\item \texttt{PNEUMONIA/person1175\_bacteria\_3122.jpeg}
\item \texttt{PNEUMONIA/person1180\_virus\_2015.jpeg}
\item \texttt{PNEUMONIA/person1181\_bacteria\_3129.jpeg}
\item \texttt{PNEUMONIA/person1183\_bacteria\_3131.jpeg}
\item \texttt{PNEUMONIA/person1183\_virus\_2018.jpeg}
\item \texttt{PNEUMONIA/person1184\_virus\_2019.jpeg}
\item \texttt{PNEUMONIA/person1187\_virus\_2023.jpeg}
\item \texttt{PNEUMONIA/person1190\_virus\_2031.jpeg}
\item \texttt{PNEUMONIA/person1191\_virus\_2032.jpeg}
\item \texttt{PNEUMONIA/person1208\_bacteria\_3160.jpeg}
\item \texttt{PNEUMONIA/person1227\_virus\_2078.jpeg}
\item \texttt{PNEUMONIA/person1244\_bacteria\_3200.jpeg}
\item \texttt{PNEUMONIA/person1281\_virus\_2204.jpeg}
\item \texttt{PNEUMONIA/person1287\_virus\_2210.jpeg}
\item \texttt{PNEUMONIA/person12\_bacteria\_48.jpeg}
\item \texttt{PNEUMONIA/person1357\_virus\_2338.jpeg}
\item \texttt{PNEUMONIA/person1358\_virus\_2339.jpeg}
\item \texttt{PNEUMONIA/person1368\_virus\_2353.jpeg}
\item \texttt{PNEUMONIA/person1405\_bacteria\_3566.jpeg}
\item \texttt{PNEUMONIA/person1405\_virus\_2408.jpeg}
\item \texttt{PNEUMONIA/person1409\_bacteria\_3583.jpeg}
\item \texttt{PNEUMONIA/person1449\_bacteria\_3746.jpeg}
\item \texttt{PNEUMONIA/person1453\_bacteria\_3770.jpeg}
\item \texttt{PNEUMONIA/person1454\_bacteria\_3782.jpeg}
\item \texttt{PNEUMONIA/person1455\_virus\_2496.jpeg}
\item \texttt{PNEUMONIA/person1459\_bacteria\_3797.jpeg}
\item \texttt{PNEUMONIA/person1461\_virus\_2510.jpeg}
\item \texttt{PNEUMONIA/person1475\_virus\_2558.jpeg}
\item \texttt{PNEUMONIA/person1483\_virus\_2574.jpeg}
\item \texttt{PNEUMONIA/person1490\_virus\_2596.jpeg}
\item \texttt{PNEUMONIA/person20\_bacteria\_69.jpeg}
\item \texttt{PNEUMONIA/person21\_bacteria\_73.jpeg}
\item \texttt{PNEUMONIA/person275\_virus\_565.jpeg}
\item \texttt{PNEUMONIA/person289\_virus\_593.jpeg}
\item \texttt{PNEUMONIA/person293\_bacteria\_1379.jpeg}
\item \texttt{PNEUMONIA/person294\_bacteria\_1386.jpeg}
\item \texttt{PNEUMONIA/person294\_virus\_606.jpeg}
\item \texttt{PNEUMONIA/person295\_bacteria\_1390.jpeg}
\item \texttt{PNEUMONIA/person295\_virus\_612.jpeg}
\item \texttt{PNEUMONIA/person299\_bacteria\_1419.jpeg}
\item \texttt{PNEUMONIA/person301\_bacteria\_1429.jpeg}
\item \texttt{PNEUMONIA/person302\_bacteria\_1430.jpeg}
\item \texttt{PNEUMONIA/person303\_bacteria\_1431.jpeg}
\item \texttt{PNEUMONIA/person311\_bacteria\_1453.jpeg}
\item \texttt{PNEUMONIA/person317\_bacteria\_1473.jpeg}
\item \texttt{PNEUMONIA/person327\_bacteria\_1509.jpeg}
\item \texttt{PNEUMONIA/person343\_bacteria\_1584.jpeg}
\item \texttt{PNEUMONIA/person344\_bacteria\_1585.jpeg}
\item \texttt{PNEUMONIA/person357\_bacteria\_1639.jpeg}
\item \texttt{PNEUMONIA/person359\_bacteria\_1642.jpeg}
\item \texttt{PNEUMONIA/person384\_bacteria\_1755.jpeg}
\item \texttt{PNEUMONIA/person38\_bacteria\_196.jpeg}
\item \texttt{PNEUMONIA/person391\_bacteria\_1782.jpeg}
\item \texttt{PNEUMONIA/person409\_bacteria\_1824.jpeg}
\item \texttt{PNEUMONIA/person410\_bacteria\_1825.jpeg}
\item \texttt{PNEUMONIA/person411\_bacteria\_1826.jpeg}
\item \texttt{PNEUMONIA/person416\_bacteria\_1840.jpeg}
\item \texttt{PNEUMONIA/person422\_bacteria\_1853.jpeg}
\item \texttt{PNEUMONIA/person423\_virus\_869.jpeg}
\item \texttt{PNEUMONIA/person432\_virus\_881.jpeg}
\item \texttt{PNEUMONIA/person444\_bacteria\_1927.jpeg}
\item \texttt{PNEUMONIA/person450\_bacteria\_1941.jpeg}
\item \texttt{PNEUMONIA/person454\_virus\_938.jpeg}
\item \texttt{PNEUMONIA/person458\_virus\_945.jpeg}
\item \texttt{PNEUMONIA/person45\_bacteria\_222.jpeg}
\item \texttt{PNEUMONIA/person46\_bacteria\_225.jpeg}
\item \texttt{PNEUMONIA/person472\_virus\_969.jpeg}
\item \texttt{PNEUMONIA/person476\_bacteria\_2026.jpeg}
\item \texttt{PNEUMONIA/person477\_bacteria\_2031.jpeg}
\item \texttt{PNEUMONIA/person488\_bacteria\_2062.jpeg}
\item \texttt{PNEUMONIA/person489\_bacteria\_2067.jpeg}
\item \texttt{PNEUMONIA/person491\_bacteria\_2071.jpeg}
\item \texttt{PNEUMONIA/person492\_bacteria\_2085.jpeg}
\item \texttt{PNEUMONIA/person493\_bacteria\_2087.jpeg}
\item \texttt{PNEUMONIA/person494\_bacteria\_2088.jpeg}
\item \texttt{PNEUMONIA/person496\_bacteria\_2095.jpeg}
\item \texttt{PNEUMONIA/person504\_bacteria\_2127.jpeg}
\item \texttt{PNEUMONIA/person504\_bacteria\_2133.jpeg}
\item \texttt{PNEUMONIA/person506\_bacteria\_2138.jpeg}
\item \texttt{PNEUMONIA/person534\_bacteria\_2251.jpeg}
\item \texttt{PNEUMONIA/person536\_virus\_1064.jpeg}
\item \texttt{PNEUMONIA/person545\_bacteria\_2290.jpeg}
\item \texttt{PNEUMONIA/person549\_bacteria\_2303.jpeg}
\item \texttt{PNEUMONIA/person575\_bacteria\_2374.jpeg}
\item \texttt{PNEUMONIA/person58\_bacteria\_278.jpeg}
\item \texttt{PNEUMONIA/person590\_bacteria\_2428.jpeg}
\item \texttt{PNEUMONIA/person610\_bacteria\_2475.jpeg}
\item \texttt{PNEUMONIA/person618\_bacteria\_2489.jpeg}
\item \texttt{PNEUMONIA/person630\_bacteria\_2512.jpeg}
\item \texttt{PNEUMONIA/person645\_bacteria\_2537.jpeg}
\item \texttt{PNEUMONIA/person652\_bacteria\_2544.jpeg}
\item \texttt{PNEUMONIA/person674\_bacteria\_2568.jpeg}
\item \texttt{PNEUMONIA/person68\_bacteria\_336.jpeg}
\item \texttt{PNEUMONIA/person703\_virus\_1300.jpeg}
\item \texttt{PNEUMONIA/person72\_bacteria\_354.jpeg}
\item \texttt{PNEUMONIA/person732\_virus\_1353.jpeg}
\item \texttt{PNEUMONIA/person734\_bacteria\_2637.jpeg}
\item \texttt{PNEUMONIA/person745\_bacteria\_2648.jpeg}
\item \texttt{PNEUMONIA/person74\_bacteria\_363.jpeg}
\item \texttt{PNEUMONIA/person781\_virus\_1412.jpeg}
\item \texttt{PNEUMONIA/person788\_virus\_1419.jpeg}
\item \texttt{PNEUMONIA/person802\_bacteria\_2708.jpeg}
\item \texttt{PNEUMONIA/person805\_bacteria\_2712.jpeg}
\item \texttt{PNEUMONIA/person813\_virus\_1449.jpeg}
\item \texttt{PNEUMONIA/person815\_bacteria\_2726.jpeg}
\item \texttt{PNEUMONIA/person833\_virus\_1469.jpeg}
\item \texttt{PNEUMONIA/person844\_virus\_1487.jpeg}
\item \texttt{PNEUMONIA/person86\_virus\_159.jpeg}
\item \texttt{PNEUMONIA/person885\_bacteria\_2809.jpeg}
\item \texttt{PNEUMONIA/person8\_bacteria\_37.jpeg}
\item \texttt{PNEUMONIA/person929\_bacteria\_2854.jpeg}
\item \texttt{PNEUMONIA/person94\_virus\_176.jpeg}
\end{itemize}

\begin{thebibliography}{9}
\bibitem{rajaraman2018} Rajaraman S.\ et al.\ Pre-trained convolutional
neural networks as feature extractors toward improved malaria parasite
detection in thin blood smear images. \emph{PeerJ} 6:e4568, 2018.
\bibitem{kermany2018} Kermany D.\ et al.\ Identifying medical diagnoses and
treatable diseases by image-based deep learning. \emph{Cell} 172(5):1122--1131, 2018.
\bibitem{kipf2017} Kipf T., Welling M.\ Semi-supervised classification with
graph convolutional networks. \emph{ICLR}, 2017.
\bibitem{northcutt2021} Northcutt C., Jiang L., Chuang I.\ Confident
learning: estimating uncertainty in dataset labels. \emph{JAIR} 70:1373--1411, 2021.
\bibitem{kingma2015} Kingma D., Ba J.\ Adam: a method for stochastic
optimization. \emph{ICLR}, 2015.
\end{thebibliography}
\end{document}
